% !TEX program = xelatex
% 2027 新领军预测训练卷 · 设计说明、真题分析与审核记录
\documentclass[cjk=noto, nomargin]{mojing-notes}
\usepackage{longtable}

\course[Design Notes]{预测训练卷设计说明}
\term{2026 年 10 月}
\instructor{Claude（依 xinlingjun-exam-expert 技能编写）}
\institute{训练资料，非官方文件}

\begin{document}

\maketitle
\tableofcontents

% =============================================================================
\lecture[number=1, date=2026 年 10 月, short=卷子的身份与依据]{这套卷子是什么，依据是什么}
% =============================================================================

\begin{keypoints}
\begin{itemize}
  \item 全套 12 份试卷：一试（数学 1、数学 2、物理），二试（语文、英语，综合六科），三试（现学现考）；每份另有参考答案及评分标准。
  \item 全部为本次编写的模拟训练题，不是官方试题，也不是历年真题。官方没有公布任何一试的题量、时长、分值，二试、三试也没有考纲；卷面安排全是训练设置。
  \item 一试依据 2027 年数学、物理考纲（2026-08-20 首发）；二试、三试只有简章中的科目名称和书院负责人的公开解释，设计带有推测成分，下文逐项说明。
\end{itemize}
\end{keypoints}

\section{官方结构}

2027 年招生简章（2026-09-28）规定的学科能力测试为三试：
\begin{itemize}
  \item \textbf{一试}：中学数学全部内容；数学 1（分析，代数中的线性代数与群论，几何与拓扑）；数学 2（概率与统计含组合，应用数学中的数值分析、数值代数、最优化）；物理（以高中核心内容为基准，超出范围的新概念由题目提供）。
  \item \textbf{二试}：语言（语文、英语）；综合，在物理加强、化学、生物、计算机、经济、社科六科中选两科。经济与社科是两个独立科目。
  \item \textbf{三试}：非固定形式的开放性内容（现学现考）。
\end{itemize}
数学基础测试（机考）是学科能力测试之前的独立环节，考纲没有给出它的细目，不在本套卷子范围内。

\section{网络资料能说明什么}

\begin{itemize}
  \item 求真书院负责人接受光明日报采访（2026-09-29，清华新闻网转载）时解释：数学考查扩展到几何与拓扑、概率统计、应用数学，强调数理并重；
        二试综合“意在考查学生在数学之外，是否怀有真实而持久的兴趣”，语言测试关注阅读、表达与交流；书院不支持针对选拔的培训。
  \item 多家媒体（北京日报、人民网、中国教育在线等）的报道只复述简章结构，没有题型、时长细节。截至 2026-10-18 未检索到 2027 年任何一场的真实考题或回忆。
  \item 往年“二试”（开放测试）的参试者回忆：先看教学视频或讲义、当场作答，涉及过狭义相对论、纽结理论等陌生专题。新版三试“现学现考”与之形式相近，但简章没有承诺沿用。
\end{itemize}
据此，二试综合各卷侧重“用高中核心知识加上题给新信息做推理”，而非大学课程的知识量；三试采用“讲义 + 由浅入深的问题 + 开放问题”的形式。这些都是训练设计，不是官方信息。

% =============================================================================
\lecture[number=2, date=2026 年 10 月, short=历年真题与经典结论]{历年一试真题也是“经典结论”吗}
% =============================================================================

\begin{keypoints}
\begin{itemize}
  \item 问题：如果一道题能叫出背后的定理名称，见过它的同学就有优势。原来的真题有没有这种性质？
  \item 结论：有，而且很普遍。按本文的判断，48 道旧一试题中约 22 道（46\%）以可以点名的经典结论为背景；2022 年 9 月一试 8 题全部如此。
  \item 但真题从不点名，常换外衣或换数域、换成特例、拆成引导性小问，并且要求证明。见过结论能帮你确定方向，却不能代替现场证明。
  \item 零试中甚至出现过原题重复（如 TACA-202208-ZERO-Q05 与 TACA-202408-ZERO-Q01），说明“见过”在历史上确实带来过直接优势。
\end{itemize}
\end{keypoints}

\section{分类方法}

读取题库中 2021—2025 年 6 份 TACA 一试共 48 道大题的完整题面，逐题归入三类：
\textbf{A 类}——题目是某个有名称的定理或经典结论的特例、变形或直接复述；
\textbf{B 类}——没有专门名称，但属于教材习题或标准技巧；
\textbf{C 类}——看不出现成的背景，结构较新颖。分类是本文的判断，不是官方信息，边界题可能有不同看法。

\section{逐题归类}

{\small
\begin{longtable}{p{2.2cm}p{4.6cm}p{5.3cm}c}
\toprule
题号 & 题目要点 & 可点名的背景 & 类 \\
\midrule
\endhead
2021.03 Q1 & $x^a\sin(1/x^3)$ 的可导性 & 振荡函数反例型习题 & B \\
2021.03 Q2 & 反常积分算子迭代 & — & C \\
2021.03 Q3 & 洛伦兹型双线性形式的反向柯西不等式 & 类时向量的反向柯西—施瓦茨（Aczél 型） & A \\
2021.03 Q4 & 无界多项式迭代的乘积增长 & — & C \\
2021.03 Q5 & 倒数自变量的微分函数方程 & — & C \\
2021.03 Q6 & $\operatorname{Aut}(\Z_2\oplus\Z_8)$ & 有限交换群自同构计算 & B \\
2021.03 Q7 & 两个椭球同时被平面三等分 & — & C \\
2021.03 Q8 & 瑞利商、锥分解、非对角元为负的矩阵正定 & Rayleigh—Ritz、Moreau 分解、M-矩阵 & A \\
2021.10 Q1 & 端点集中的积分极限 & 标准技巧 & B \\
2021.10 Q2 & $x^3+y^3+z^3+txyz$ 分解为一次因式 & Hesse 三次曲线束 & A \\
2021.10 Q3 & 对数误差与二重黎曼和 & 标准技巧 & B \\
2021.10 Q4 & 零对角符号矩阵的秩 & — & C \\
2021.10 Q5 & 矩阵多项式映射的满射性 & — & C \\
2021.10 Q6 & $f''+2021f\ge0$ 的正值区间长度 & Sturm 比较定理 & A \\
2021.10 Q7 & 循环群直积商的元素最大阶 & 有限交换群结构 & B \\
2021.10 Q8 & $\ln(1+x)$ 迭代的首阶渐近 & 经典迭代渐近（同 $\sin$ 迭代） & B \\
2022.09 Q1 & 半圆权函数的矩 & Catalan 数（半圆律的矩） & A \\
2022.09 Q2 & $X\mapsto AX+XA$ 的可逆性 & Sylvester 方程定理 & A \\
2022.09 Q3 & 不定二次型的全迷向子空间 & Witt 指数 & A \\
2022.09 Q4 & $V^2\le S_1S_2S_3$ & Loomis—Whitney 不等式 & A \\
2022.09 Q5 & $\mathrm{SL}_2(\Z)$ 的生成元 & 模群由 $S,T$ 生成 & A \\
2022.09 Q6 & 有限可消去半群是群 & 经典定理 & A \\
2022.09 Q7 & $\cosh x\le\ee^{x^2/2}$ 与加权指数矩 & Hoeffding 引理（变体） & A \\
2022.09 Q8 & 双曲整数变换下不变的周期连续函数 & 双曲环面自同构的遍历性（猫映射） & A \\
2023.10 Q1 & 素数计数函数的积分 & — & C \\
2023.10 Q2 & 奇函数复合交换子的消失阶 & — & C \\
2023.10 Q3 & $AB+A=BA+B\Rightarrow(A-B)^n=0$ & 交换子幂零（Jacobson 引理型） & A \\
2023.10 Q4 & 三维格点级数的收敛阈值 & 标准技巧 & B \\
2023.10 Q5 & 两个 Jordan 块的中心化子维数 & Frobenius 中心化子维数公式 & A \\
2023.10 Q6 & $X\mapsto AXA\T$ 的迹与行列式 & 对称平方 & B \\
2023.10 Q7 & 标准表示特征标的平方和 & 特征标正交关系 & A \\
2023.10 Q8 & 等范数向量的统一常数 & — & C \\
2024.11 Q1 & 指数权重下的端点集中积分 & 标准技巧 & B \\
2024.11 Q2 & 右矩阵模同态的维数 & 模同态空间 & B \\
2024.11 Q3 & 倍周期函数空间上的微分算子 & — & C \\
2024.11 Q4 & 有理正交变换分解为反射之积 & Cartan—Dieudonné 定理（有理数域） & A \\
2024.11 Q5 & 传递交换置换群 & 传递交换作用必为正则作用 & A \\
2024.11 Q6 & 椭圆型方程零边界解唯一 & 极值原理 / 能量估计 & A \\
2024.11 Q7 & 二进制序列测试族刻画连续性 & — & C \\
2024.11 Q8 & 非阿基米德绝对值的分类 & Ostrowski 定理 & A \\
2025.11 Q1 & 下确界的渐近 & — & C \\
2025.11 Q2 & $S(A^k)\ge S(A)^k/n^{k-1}$ & Blakley—Roy 型不等式 & A \\
2025.11 Q3 & 镜面对称换位生成的置换群 & 超八面体群（B 型 Weyl 群） & A \\
2025.11 Q4 & 非线性常微分方程解的渐近 & 与 2021.10 Q8 同类 & B \\
2025.11 Q5 & 自由单轨作用下的有向图分类 & Cayley 图 & B \\
2025.11 Q6 & 顺序随机删除流程 & — & C \\
2025.11 Q7 & 模 2 可逆块矩阵计数 & $\abs{\mathrm{GL}_2(\F_4)}$ & B \\
2025.11 Q8 & 主值积分的连续性与可积性 & Hilbert 变换 & A \\
\bottomrule
\end{longtable}
}

合计：A 类 22 题，B 类 13 题，C 类 13 题。

\section{真题怎样处理“经典背景”}

\begin{enumerate}
  \item \textbf{不点名}：没有一道题写出定理名称，考生须自己认出结构。
  \item \textbf{换外衣或换数域}：Cartan—Dieudonné 放到有理数域上；Loomis—Whitney 只对凸多面体；Hoeffding 引理改写成确定性的加权不等式。
  \item \textbf{拆成引导小问}：2024 年第 8 题把 Ostrowski 定理拆成三步，每一步都要证明。
  \item \textbf{要证明而不是要答案}：大部分 A 类题是证明题，背出结论只能确定目标。
\end{enumerate}

\section{对本套卷子的影响}

第一版预测卷在题面上直接写出了“$P(G)\le\frac58$”“7 种颜色”和 Kantorovich 常数，等于把答案交给了见过的人。按真题的处理方式，本版作了以下修改：
\begin{itemize}
  \item \textbf{交换概率}（数学 1 第 5 题）：不再给出 $\frac58$，改问“$\abs G$ 的素因子都不小于 $p$”时的最大值 $\frac{p^2+p-1}{p^3}$ 及取等条件；
        第二轮审核后又加入“取到最大值时换位子群是 $p$ 阶中心子群”（要用换位子的双线性）。第三轮审核指出原拟的“哪些阶可以取到”
        由 $p^3k(G)=(p^2+p-1)\abs G$ 一行算术即得，已据此改成现在的结构性问题。
  \item \textbf{环面染色}（数学 1 第 6 题）：不再给出“7”，先问 $n$ 顶点可嵌入图的最大边数 $e(n)=3n$（需要构造一族三角剖分）；
        第二轮又加入二部图的版本 $b(n)=2n$（$n\ge8$），构造要从三角格换成正方形格。
  \item \textbf{立方体图生成树}（数学 1 第 4 题）：一阶极限 $\tau_n^{1/2^n}\sim n$ 容易猜，第二轮改问 $\lim(\tau_n^{1/2^n}-n)=-\frac12$，
        需要四阶矩控制的 Taylor 估计。
  \item \textbf{对数凹性}（数学 1 第 3 题）：第一版用正单形向量，审核者指出“迹恒为 1 + 均值不等式”一行即得答案；第二轮换成迹不恒定、
        也不对称的向量组 $e_1,\dots,e_n,\,e_1+\dots+e_n$，并要求证明最大值点唯一（要用 (2) 的严格性）。
  \item \textbf{德国坦克问题}（数学 2 第 3 题）：加入“编号起点也未知”，需要用间隔的可交换性求极差的方差，答案 $\frac{2n}{n-1}$ 不在常见资料中。
  \item \textbf{最速下降}（数学 2 第 6 题）：改问“最小常数”；第二轮把二维情形推广到任意初始点，证明每步比值不变并求出它，
        取等的初始点作为推论。
  \item \textbf{椭圆内切圆盘}（数学 1 第 1 题，原创）：第二轮加入“极大圆盘圆心的集合”（椭圆的中轴），分值由 20 提到 25。
\end{itemize}
按同样的标准自评：一试 12 道数学题中，A 类约 8 道、B 类约 3 道、C 类 1 道，A 类比例仍高于真题的 46\%，但每道 A 类题都带有背诵结论无法直接解决的变式，且第二轮审核点名的“背诵优势高”的小问都已加入新的一问或换掉。
继续优化时可以再换入 C 类原创题。

% =============================================================================
\lecture[number=3, date=2026 年 10 月, short=一试设计]{一试三卷的设计说明}
% =============================================================================

\begin{keypoints}
\begin{itemize}
  \item 数学 1、数学 2 各 6 题、150 分、建议 3 小时；物理 5 题、100 分、建议 2 小时（训练设置）。
  \item 按老师的命题要求：题面轻、思维量大；较难题有 2—3 个相互依赖的关键观察并融合多个板块；计算量以 2023 年 10 月一试第 2 题为定性参照。
  \item 老师建议该同学“主要学概率论数理统计以及应用数学”，且历史题库中统计、数值代数、物理几乎为空，因此数学 2 与物理按新考纲补足。
\end{itemize}
\end{keypoints}

\section{逐题要点}

{\small
\begin{longtable}{p{1.5cm}p{3.2cm}p{5.9cm}p{2.5cm}}
\toprule
题 & 考纲方向 & 关键观察 & 背景 \\
\midrule
\endhead
数 1-1 & 解析几何、曲线曲率 & 相切推出圆心在法线上；整体障碍（对称点）与局部障碍（曲率）；嵌套圆盘刻画极大圆盘 & 原创（中轴） \\
数 1-2 & 傅里叶级数、不等式 & 分段柯西不等式；锯齿核表示公式或 Parseval；用光滑函数逼近说明最优 & 最优 Sobolev 常数 \\
数 1-3 & 多元积分、谱定理、群对称 & 正交换元算 Gauss 积分；Hölder 及其取等给严格对数凹性；对称群平均 + 严格凹性得唯一最优权重 & Ky Fan 不等式、D-最优设计 \\
数 1-4 & 线性代数、特征标、组合、概率 & 用 $\{0,1\}^n$ 的特征标对角化；矩阵树定理；对数的 Taylor 展开 + 四阶矩估计余项 & 立方体图生成树 \\
数 1-5 & 群作用、计数 & 交换对 $=\abs G\,k(G)$；$G/Z$ 循环则交换；素因子下界同时抬高两处估计；换位子双线性给出 $\abs{G'}=p$ & 5/8 定理的推广 \\
数 1-6 & 组合拓扑、图论 & 区域边界含圈推出 $E\le3V$ 与二部图的 $E\le2V$；三角格、正方形格模子格构造；贪心染色 & Heawood \\
数 2-1 & 随机游走、谱 & 交换求和化为二次型；三对角矩阵与余弦试探解；两端边界条件 & $\min(i,j)$ 矩阵 \\
数 2-2 & 分支过程、渐近 & 首步分析得母函数迭代；倒数变换；二阶修正中的对数项 & Kolmogorov 估计 \\
数 2-3 & 抽样、估计 & 间隔可交换（隔板法）；极差与最大值的方差；无偏修正 & 德国坦克问题 \\
数 2-4 & 假设检验、群作用 & 不变性取平均；群使参照集合不变；含并列的计数引理；非群反例 & 置换检验有效性 \\
数 2-5 & 数值积分、傅里叶 & 等比级数给傅里叶系数；混叠公式；Richardson 外推的前提失效 & 周期梯形公式 \\
数 2-6 & 最优化、谱分解 & 用 $g\T A^{-1}g$ 表示目标值；矩阵不等式 + 均值不等式；二维中相邻梯度正交使比值不变 & Kantorovich \\
物 -1 & 热学、振动、误差 & 平衡压强计入重力；绝热线性化；两个时间尺度；壁面层使 $\gamma$ 偏小 & Rüchardt 实验 \\
物 -2 & 电磁感应 & 动量—电量关系；电容相当于质量 $B^2L^2C$；效率上限 $\frac14$ & 经典模型 \\
物 -3 & 近代物理、数据处理 & 最小二乘；系统误差只移截距；最优实验设计与检验线性 & 光电效应 \\
物 -4 & 万有引力、能量 & 动能 $=-E$ 推出阻力使卫星加速；引力切向分量为 $2f$ & 卫星阻力佯谬 \\
物 -5 & 干涉、拍 & 光程差与条纹计数；双线“拍”与可见度；相干长度 & 迈克耳孙干涉仪 \\
\bottomrule
\end{longtable}
}

\section{考纲覆盖与缺口}

一试覆盖了考纲中的：解析几何、曲率；傅里叶级数、多元积分；谱定理、二次型；群作用与计数、线性群；欧拉数、三角剖分、图论、染色；
随机游走、简单随机过程、条件概率；抽样、极大似然、矩估计、无偏估计、两类错误、p 值、对称性检验；数值积分；迭代法、Cholesky 分解；梯度下降；
热学、电磁感应、万有引力、干涉、光电效应、实验数据处理、数量级估算。

未覆盖、建议另行补练：西罗定理、自由群与关系、舒尔引理、若尔当标准形、张量积、实数轴的紧致与完备、凸包；
线性规划、非线性方程求根、插值与逼近、特征值的数值计算；中心极限定理；几何光学、相对论（三试模块 B 有所涉及）。

% =============================================================================
\lecture[number=4, date=2026 年 10 月, short=二试与三试]{二试、三试的设计依据与不确定性}
% =============================================================================

\section{二试语言}

官方只列出“语言（语文、英语）”。语文卷按“现代文阅读、文言文阅读、语言运用、写作”组织，阅读材料为原创议论文，古文选刘徽割圆术注与《墨经》几何定义，
写作题围绕“探索与继承”。英语卷侧重学术阅读（原创的两篇科普议论文）、英汉互译与论证写作，这与 2027 年面试含英文面试相呼应。题型与分值都是训练设置。

\section{二试综合六科}

六科都按“高中核心内容 + 题给新信息 + 定量或证据推理”编写，每卷 4 题、100 分、建议 1.5 小时：
物理加强（傅科摆、正交电磁场、可逆热机与熵、量纲分析与水波）；化学（气相平衡、缓冲溶液、电化学、动力学）；
生物（哈迪—温伯格平衡、连锁与交换、DNA 复制实验推理、逻辑斯谛模型与捕捞）；计算机（比较次数下界、动态规划、BFS 与 Dijkstra、汉明码）；
经济（税负归宿、重复博弈、贷款计算、比较优势）；社科（投票悖论、辛普森悖论、抽样偏差、相关与因果）。

\textbf{最大的不确定性}：官方没有任何一科的考纲，也没有说明是否上机（计算机）、是否以材料题为主（社科）。若真实考试偏重知识面，本卷对知识量的覆盖是不足的。

\section{三试现学现考}

两个模块各 50 分：数学模块讲四元数与空间旋转，物理模块讲快度与相对论速度合成。每个模块先给一页讲义，再按“复述定义—直接推论—与已知工具联系—变式—开放问题”出题，
对应 xinlingjun-exam-expert 技能中“现学现考训练”的检查顺序。真实三试的形式（讲座、视频、讲义、面谈）未知。

% =============================================================================
\lecture[number=5, date=2026 年 10 月, short=审核与使用]{独立审核记录与使用建议}
% =============================================================================

\section{第一版的审核（2026-10-18）}

第一版（9 道数学 + 3 道物理）由三位读者分组盲做，再比对命题者解答。全部小问答案一致；据意见修改了第 10 题（原阻尼小问没有给出模型，且与“无摩擦”设定下的主要阻尼机制矛盾，已删去重写），
补全了环面区域边界的论证，并在第一问加上“正则”条件等。详细记录见第一版的说明文件。

\section{本版的审核}

（审核进行中，完成后补入。）


\section{使用建议}

\begin{enumerate}
  \item \textbf{按场次限时}：一试三卷各按建议用时独立完成；二试语文、英语各 2 小时，综合任选两科各 1.5 小时；三试 2.5 小时（含阅读）。
  \item \textbf{先看速查，再读详解}：对每道没做出的题标出卡点类型——K 知识、R 表示、D 结构发现、P 证明闭合、C 计算执行。
  \item \textbf{按卡点补练}：K 补对应考纲条目；D 只看一个关键观察后合上答案重做；P 重写完整证明；C 做同类短计算。
  \item \textbf{对 A 类题}：先不看背景名称独立做；做完后再查它的经典来源，比较题目的变式改动了什么。
  \item \textbf{二试、三试}：重点不是背知识，而是读懂新材料并用它推理；建议把讲义题当作“限时自学”练习。
\end{enumerate}

\end{document}
